\subsection{Hamilton's Method}


\content{
        \begin{definition} For \textbf{Hamilton's Method}, we perform the
        following steps: 
        \begin{itemize}
        \item Calculate the \textbf{divisor}: (total population divided by total
        number of representatives) 
        \item Calculate the \textbf{quota} for each state: (state's population
        divided by divisor)
        \item Calculate \textbf{lower quotas}
        \item Assign any remaining representatives to those states whose decimal
        parts of the quota was largest. 
        \end{itemize}
        \end{definition}
}{% presentation
}{% nopresentation
}{% comments
lower quotas are rounded down. 
}


\content{
\begin{example} The state of Delaware has three counties: Kent, New Castle, 
and Sussex. The Delaware state House of Representatives has 41 members. If 
Delaware wants to divide this representation along county lines,
let’s use Hamilton’s method to 
apportion them. The populations of the counties are as follows (from the 
2010 Census): Kent: 162,310, New Castle: 538,479, and Sussex: 197,145.
\end{example}
}{% presentation
\vspace{3in}
}{% nopresentation
\vspace{3in}
}{% comments
total pop is 897,934.
}

\content{
        \begin{example} Use Hamilton's method to apportion the 75 seats of Rhode
        Island's House of Representatives amoung it's five counties. \\
        \begin{tabular}{lr}
        County & Population \\
        \hline
        Brisol & 49,875 \\
        Kent & 166,158\\
        Newport & 82,888 \\
        Providence & 626,667\\
        Washington & 126,979 \\
        \hline
        Total & 1,052,567
        \end{tabular}
        \end{example}
}{% presentation
\vspace{3in}
}{% nopresentation
\vspace{3in}
}{% comments
in order, final answer is 3, 12, 6, 45, 9. 
}

\content{
        \begin{example}
        Three peaople invest in a treasure dive, each investing the amount
        listed below. The dive results in 36 gold coind. Apportion those coins
        to the investors. \\
        \begin{tabular}{lr}
        Alice & \$7,600 \\
        Ben: & \$5,900 \\
        Calos: & \$1,400
        \end{tabular}
        \end{example}
}{% presentation
\vspace{3in}
}{% nopresentation
\vspace{3in}
}{% comments
}

\content{
        \begin{example} A state with five counties has 50 seats in their
        legislature.
        Using Hamilton’s method, apportion the seats based on the 2000 census, 
        then again using the 2010 census. \\
        \begin{tabular}{lll}
        County & 2000 Population & 2010 Population \\
        \hline
        Jefferson & 60,000 & 60,000 \\
        Clay & 31,200 & 31,200 \\
        Madison & 69,200 & 72,400 \\
        Jackson & 81,600 & 81,600 \\
        Franklin & 118,000 & 118,400\\
        \hline
        \end{tabular}
        \end{example}

}{% presentation
\vspace{4in}
}{% nopresentation
\vspace{4in}
}{% comments
talk about what goes wrong here, this is called the population paradox. 
}
